% Root derivation; prospective B-owned label, not included at creation.
\paragraph{The complete Greenberg functor on its stated test rings.}
The complete construction also has a precise modern theorem.  We specify
the test rings because this theorem does not interchange a finite
partition colimit with an inverse limit.

\begin{theorem}[Known, with explicit derivation: complete Greenberg points]\label{B:complete-Greenberg-points}
Let $R$ be a complete discrete valuation ring with perfect residue field
$k$ and uniformizer $\pi$.  In equal characteristic fix an isomorphism
$R\simeq k[[t]]$ carrying $\pi$ to $t$.  In mixed characteristic use the
canonical Cohen inclusion $W(k)\hookrightarrow R$.  Put $R_n=R/(\pi^n)$
and let $\mathcal R_n$ be its finite Greenberg ring scheme over $k$.
For a $k$-algebra $A$, assume that $A$ is perfect in mixed characteristic,
and define
\[
 S(A)=\varprojlim_n\mathcal R_n(A).
\]
For every $R$-scheme $X$, without a finite-type or separatedness assumption,
the scheme
\[
 \operatorname{Gr}_R(X)=
 \varprojlim_n\operatorname{Gr}_{R_n}(X\times_R R_n)
\]
exists.  There are natural bijections
\[
 \operatorname{Gr}_R(X)(A)
 \simeq\varprojlim_n X\bigl(\mathcal R_n(A)\bigr)
 \simeq X\bigl(S(A)\bigr),
\]
where all points on the right are morphisms over $R$.  Here
\[
 S(A)\simeq
 \begin{cases}
 A[[t]],&R\simeq k[[t]],\\
 R\otimes_{W(k)}W(A),&\operatorname{char}R=0,\ \operatorname{char}k>0.
 \end{cases}
\]
The construction preserves products and group objects.  It is generally
not of finite type: for $X=\mathbb A^1_R$ its underlying scheme is
$\operatorname{Spec}k[x_0,x_1,\ldots]$.
\end{theorem}

\begin{proof}[Derivation from the opened primary theorems]
The finite ring schemes and their transition maps are those of
\cite[Sections~4,6--7]{BBertapelleGonzalezAviles2022}.
The transition morphisms in the displayed system of Greenberg schemes
are affine; Proposition--Definition~17.5 and its proof give the inverse
limit and its functor of points
\cite[pp.~93--94; see also (18.1)]{BBertapelleGonzalezAviles2022}.
This proves the first bijection using the finite-level point formula
(7.16).  No exchange with a filtered colimit has been made.

In equal characteristic, the ring formula and
$S(A)/(t^n)=\mathcal R_n(A)$ are immediate from
$\mathcal R_n(A)=A[t]/(t^n)$.  In mixed characteristic,
\cite[Remark~6.6(b), (6.7)]{BBertapelleGonzalezAviles2022}
identifies $S(A)$ with $R\otimes_{W(k)}W(A)$ and its quotient by
$\pi^n$ with $\mathcal R_n(A)$.  We check completeness, which is the
essential hypothesis of the next theorem.  Write
$R=W(k)[T]/(f)$ for an Eisenstein polynomial $f$ of degree $e$,
using the chosen uniformizer.  Then
$S(A)=W(A)[T]/(f)$ is finite free over the $p$-adically complete ring
$W(A)$; here perfection of $A$ gives
$W(A)/(p^n)=W_n(A)$.  It is consequently $p$-adically complete.
The Eisenstein relation gives $\pi^e=p u$ with $u$ a unit:
the relevant factor has a unit constant term and all its other terms
are divisible by $\pi$, which is topologically nilpotent since
$\pi^e\in pS(A)$.  Thus the $p$-adic and $\pi$-adic topologies are
cofinal.  We obtain in both cases
\[
 S(A)\simeq\varprojlim_n S(A)/(\pi^n),
 \qquad S(A)/(\pi^n)\simeq\mathcal R_n(A).
\]

Bhatt's Theorem~4.1 gives
$X(S(A))\simeq\varprojlim_n X(S(A)/(\pi^n))$ for qcqs algebraic
spaces; his Remark~4.6 extends it to every scheme
\cite[Section~4, pp.~16--17]{ROOT_BhattAlgebraization2014}.
Apply that theorem to the underlying scheme of $X$ and the complete
ring $S(A)$.  The bijection respects the condition of being over $R$:
if the compatible maps at every level are over $R$, the structural map
$R\to S(A)$ obtained after algebraization agrees with the prescribed
one modulo every $\pi^n$, hence agrees with it because $S(A)$ is
separated.  This proves the second bijection, including its relative
typing.  This is the corrected justification of
\cite[Proposition~18.2]{BBertapelleGonzalezAviles2022}.

The finite adjunctions preserve products, as do inverse limits; applying
these identities to multiplication, inverse, identity and their diagrams
gives the assertion about group objects.  Finally, the underlying scheme
of $\mathcal R_n$ is affine $n$-space with the compatible coordinate
projections of the construction, so its inverse limit is the infinite
affine scheme displayed above, as also recorded in
\cite[Section~6 and the paragraph after (18.1)]{BBertapelleGonzalezAviles2022}.
\end{proof}
